\section{引言}
\label{sec:intro}

水产养殖是我国重要的农业产业，而病害是制约养殖效益与食品安全的关键因素。养殖过程中，牛蛙、鱼类等水产生物一旦发病，养殖户往往只能提供\textbf{一张现场照片和一句口头主诉}（例如"今天蛙不吃料了"），然后期待一个明确的诊断与处理方案。然而，从这样有限的输入中得出可靠诊断，本质上是一个信息严重不足的推理过程：照片只能提供表观病灶，主诉只能提供行为异常，而\textbf{水温、死亡率、粪便、是否剖检、养殖密度}等关键临床字段完全缺失。一个合格的诊断助手必须在给出结论之前，主动向养殖户追问这些缺失字段。

更复杂的是，水产诊断是一个\textbf{高风险、强合规}的任务。误诊和乱开药不仅直接造成经济损失，还涉及禁用药与处方药的食品安全红线；而"看不清就送检"是兽医执业的基本要求。这意味着一个可用的诊断模型不仅要"会追问、会诊断"，还必须"该拒答时明确拒答""全程不给出具体用药剂量"。这三个行为——\textbf{追问、拒答、合规}——共同构成了专业诊断的安全底线。

然而，我们发现现有视觉语言模型（VLM）在这三个行为上存在系统性的失败。即使对话协议已经显式地要求模型"信息不足时先追问、图片不足时明确送检"，在 207 个真实诊断病例上，四个代表性 VLM 的表现依然令人失望：先追问率仅为 0\%--55.6\%，该拒答时的拒答准确率仅为 55.6\%--94.4\%，而合规违规率最高达到 \textbf{41.5\%}（模型直接输出药物与剂量）。更值得注意的是一种"能力互补但同样危险"的模式：诊断准确率最高的 Qwen3-VL-8B（83.1\%）恰恰是违规率最高的（41.5\%）；而最安全的 Qwen3.8-27B（违规 1.9\%）诊断准确率却只有 72.0\%。\textbf{这表明"诊断准"与"守安全"在当前 VLM 中难以兼得，且光靠提示词无法弥合这一鸿沟。}

现有工作尚未正面解决这一组合问题。多模态问答中的澄清研究（如 ACL 2025 的 \textit{Teaching VLMs to Ask}~\cite{jian2025teaching}）聚焦于"问题歧义"的消解，而非"诊断信息不足"的追问；医疗对话系统（如 MedDG~\cite{liu2020meddg}）关注文本槽位填充，缺少多模态输入与合规约束；VLM 强化后训练（如 Patho-R1~\cite{zhang2026pathor1}、Reason-RFT~\cite{tan2025reasonrft}）致力于提升推理能力与效率，却没有把"何时该追问、何时该拒答"作为对齐目标。更重要的是，这些工作都没有提供一个能同时度量\textbf{追问、拒答、合规与病灶识别}的评测基准，使得"专业诊断中的主动问询与安全拒答"这一能力长期缺乏可复现的度量。

为填补这一空白，本文构建 \textbf{AquaDiag}，首个面向低资源水产的多模态诊断对话基准。AquaDiag 将诊断任务显式解耦为两个子任务：\textbf{病灶/体征识别}（确定性任务——图像中可见的病灶类型与部位，可由受控本体直接标注与评测）与\textbf{病种诊断}（不确定性任务——给定图像与多轮问诊，模型须先追问关键字段、信息足够才下诊断、无法确诊则拒答并建议送检、全程遵守合规红线）。基准包含 207 个病例、覆盖 5 个病种，提供黄金追问序列、受控病灶本体（13 种病灶类型 $\times$ 10 个部位）与 6+1 项指标。在此基础上，我们提出轻量方法 \textbf{ProDiag}，以"视觉显式化 + 可审计规则表 + 红线硬过滤"替代自由生成，在不进行任何训练的情况下，即在诊断准确率（100\%）、合规违规率（0\%）、先追问率（100\%）、必问字段覆盖率（94.2\%）与拒答准确率（100\%）上全面超越三个基线 VLM，证明了\textbf{可审计的领域规则结构能够同时修复追问、合规与拒答三类行为缺陷}。

本文的主要贡献如下：
\begin{itemize}
    \item \textbf{AquaDiag 基准}：首个面向低资源水产的多模态诊断对话基准，将任务解耦为"病灶/体征识别（确定性）"与"病种诊断（不确定性，含合规红线与拒答）"，提供 207 病例、6+1 指标、受控病灶本体与黄金追问序列。
    \item \textbf{系统性发现}：揭示三个 VLM 在"该追问不追问、该拒答不拒答、违规开药"上的系统性失败，并实证"病灶识别 $\neq$ 病种诊断"（专家标注显示 53\% 的真实病灶图无法仅凭图像确定病种，爱德华氏菌病 6/6 图均需送检）。
    \item \textbf{ProDiag 方法}：以"视觉显式化 + 规则表 + 红线过滤"的轻量可审计方法，不训练即实现诊断 100\% + 合规 0\% + 追问 100\% + 拒答 100\%，为低资源专业诊断提供了一条可行且可解释的路线。
\end{itemize}
